\chapter{Quotient Groups}

\section{Cosets and Indices}

\begin{example}
  Write \[
    S_3 = \{\id, \rho, \rho^2, \mu, \rho\mu, \rho^2\mu\},
    \quad\text{where } \rho = (123), \mu = (12).
  \]
  Let \(H = \<\mu\>\), then the left cosets of \(H\) in \(S_3\) are \[
    H = \{\id, \mu\}, \quad
    \rho H = \{\rho, \rho\mu\}, \quad
    \rho^2 H = \{\rho^2, \rho^2\mu\}.
  \]
  The right cosets of \(H\) in \(S_3\) are \[
    H = \{\id, \mu\}, \quad
    H\rho = \{\rho, \mu\rho\} = \{\rho, \rho^2\mu\}, \quad
    H\rho^2 = \{\rho^2, \mu\rho^2\} = \{\rho^2, \rho\mu\}.
  \]
\end{example}

The following theorem allows us to use number theory to study finite groups.

\begin{theorem}[Lagrange]
  Let \(G\) be a finite group and \(H\) a subgroup of \(G\), then \(|G| = [G:H]|H|\). Consequently, \(|H|\) divides \(|G|\).
\end{theorem}
\begin{proof}
  It suffices to show every cosets of \(H\) has the same size as \(H\). Consider \[
    \f:H\to aH, \quad
    \f(h) = ah.
  \]
  We check bijectivity.
  \begin{itemize}
    \item Suppose \(f(h_1) = f(h_2)\), then \[
      h_1 = a^{-1}ah_1 = a^{-1}f(h_1) = a^{-1}f(h_2) = a^{-1}ah_2 = h_2.
    \]
    \item Pick \(ah\in aH\), then\[
      \f(a^{-1}ah) = a(a^{-1}h) = ah. \qedhere
    \]
  \end{itemize}
\end{proof}

\begin{corollary}
  Let \(G\) be a finite group and \(a\in G\). Then \(|a|\) divides \(|G|\), and \(a^{|G|} = e\).
\end{corollary}
\begin{proof}
  Consider the subgroup \(\<a\>\), then \[
    |a| = |\<a\>| \Bigm| |G|.
  \]
  In particular, \(|a| = k|G|\) for some \(k\), then \[
    a^{|G|} = a^{k|a|} = (a^{|a|})^k = e^k = e. \qedhere
  \]
\end{proof}

\begin{corollary}
  Every group of prime order is cyclic.
\end{corollary}
\begin{proof}
  Let \(G\) be a group of prime order \(p\). Choose \(a\in G\setminus\{e\}\), then \(|a| > 1\) and \(|a|\) divides \(p\), so \(|a| = p\). Hence \(G = \<a\>\).
\end{proof}

\begin{exercise}
  Let \(G\) be a group of order \(p\), then any group homomorphism \(\f:G\to G'\) is either injective or trivial.
\end{exercise}
\begin{proof}
  Consider the subgroup \(\ker\f\), then \[
    |\ker\f| \Bigm| |G| = p,
  \] so \(|\ker\f| = 1\) or \(p\), which forces \(\ker\f = \{e\}\) or \(\ker\f = G\).
\end{proof}

\begin{exercise}
  Let \(G\) be an abelian group of order \(2n\) where \(n\) is odd. Then \(G\) contains precisely one element of order \(2\).
\end{exercise}
\begin{proof}
  Choose \(a\in G\setminus\{e\}\).
  \begin{itemize}
    \item If \(|a|\) is odd, then \(|a|\) divides \(n\), then
  \end{itemize}
\end{proof}

\section{Normal Subgroups}

Given \(H < G\), it is natural to ask if one can define a group structure on the set of left cosets \(\{aH | a\in G\}\) by \[
  (aH)(bH) = (ab)H
\]

It turns out that this is not well-defined unless \(H\) satisfies a condition.
\begin{definition}
  A subgroup \(H < G\) is \emph{normal} if \[
    aH = Ha \text{ for all } a\in G
  \]
  We denote a normal subgroup by \(H \lhd G\).
\end{definition}

\begin{proposition}
  Let \(H\) be a subgroup of \(G\). Then the following are equivalent:
  \begin{enumerate}
    \item \(H\) is normal in \(G\).
    \item \(aH \subset Ha\) for all \(a\in G\).
    \item \(aHa^{-1} \subset H\) for all \(a\in G\).
    \item \(aHa^{-1} = H\) for all \(a\in G\).
  \end{enumerate}
\end{proposition}
\begin{proof}
  Trivial.
\end{proof}

\begin{proposition}
  If \(G\) is an abelian group, then every subgroup \(H < G\) is normal.
\end{proposition}

\begin{example}
  \(n\Z\lhd\Z\) and \(\Z\lhd\R\).
\end{example}

\begin{example}
  \(\Z\lhd\Q\lhd\R\lhd\C\).
\end{example}

\begin{theorem}
  Let \(H\) be a subgroup of \(G\), then the operation on left cosets \[
    (aH)(bH) = (ab)H
  \] is well-defined if and only if \(H\) is normal in \(G\).
\end{theorem}
\begin{proof}
  \((\Rightarrow)\) Suppose the operation is well-defined. Let \(a\in G\), \(h\in H\), then \[
    aha^{-1}H = (ahH)(a^{-1}H)
    = (aH)(a^{-1}H) = H,
  \] so \(aha^{-1}\in H\) for all \(h\), hence \(aHa^{-1}\subset H\).

  \((\Leftarrow)\) Suppose \(H\) is normal. Let \(a,a',b,b'\in G\), suppose \(aH = a'H\) and \(bH = b'H\). 
  Then \(a' = ah\) and \(b' = bk\) for some \(h,k\in H\).
  By normality, we have \(b^{-1}h = h'b^{-1}\) for some \(h'\in H\).
  Note that \begin{align*}
    (a'b')^{-1}(ab)
    = b'^{-1}a'^{-1}ab
    = (bk)^{-1}(ah)^{-1}ab
    = k^{-1}(b^{-1}h^{-1})b
    = k^{-1}(h'b^{-1})b
    = k^{-1}h' \in H.
  \end{align*}
  Then \((a'b')^{-1}(ab)H\cap H\neq \emptyset\), which forces \((ab)H = (a'b')H\).
\end{proof}

\section{Quotient Groups}

\begin{corollary}
  Let \(H\) be a subgroup of \(G\), then \emph{quotient group} of \(G\) modulo \(H\) \[
    G/H = \{aH \mid a\in G\}
  \]
  is a group under the operation \((aH)(bH) = (ab)H\) if and only if \(H\) is normal in \(G\).
\end{corollary}

\begin{example}
  Let \(G\) be a group, then \(\{e\}\lhd G\) and \(G\lhd G\). Moreover, \[
    G/\{e\} \cong G \quad\text{and}\quad G/G \cong \{e\}.
  \]
\end{example}

\begin{example}
  \(A_n\lhd S_n\) for any \(n > 0\).
\end{example}

\begin{example}
  Let \(a = e^{i\frac{2\pi}{n}}\in D_n\), then \(\<a\>\lhd D_n\).
\end{example}
\begin{proposition}
  If \(H < G\) has \([G:H] = 2\), then \(H\lhd G\).
\end{proposition}
\begin{proof}
  Let \(a\in G\setminus H\), then \(aH\) and \(Ha\) are cosets of \(H\) in \(G\) not equal to \(H\). Since \([G:H] = 2\), we must have \(aH = Ha\) for all \(a\in G\setminus H\). Hence \(H\) is normal in \(G\).
\end{proof}

\begin{example}
  \(SO_n(\R)\lhd O_n(\R)\) for any \(n > 0\).
\end{example}
\begin{proof}
  Let \(A\in O_n(\R)\), then we have \begin{align*}
    \det(A^TA) = I_n \implies \det(A)^2 = 1 \implies \det(A) = \pm 1.
  \end{align*}
  Put \(D = \diag(-1,1,\ldots,1)\in O_n(\R)\), then \(\det(D) = -1\), so \[
    O_n(\R) = SO_n(\R)\sqcup D(SO_n(\R)),
  \] showing that \([O_n(\R):SO_n(\R)] = 2\). Hence \(SO_n(\R)\lhd O_n(\R)\).
\end{proof}

\begin{example}
  \(SL_n(\R)\lhd GL_n(\R)\) for any \(n > 0\).
\end{example}
\begin{proof}
  Let \(Q\in SL_n(\R)\), \(A\in\GL_n(\R)\), then \[
    \det(Q^{-1}AQ) = \det(Q^{-1})\det(A)\det(Q) = \frac{1}{\det(Q)}\det(A)\det(Q) = \det(A) = 1,
  \] so \(Q^{-1}AQ\in SL_n(\R)\). Hence \(SL_n(\R)\lhd GL_n(\R)\).
\end{proof}

\begin{example}
  \(GL_n(\R) \ntrianglelefteq GL_{n+1}(\R)\) for any \(n > 0\).
\end{example}
\begin{proof}
  Recall that \[
    GL_n(\R) \cong \iota(GL_n(\R)) < GL_{n+1}(\R),
  \] where \(\iota\) is the natural embedding.
  Consider \[
    A = \iota\begin{pmatrix}
      2 & \\
        & I_{n-1}
    \end{pmatrix} = \begin{pmatrix}
      2 &  & \\
        & I_{n-1} & \\
        &  & 1
    \end{pmatrix} \in GL_{n+1}(\R), \quad
    X = \begin{pmatrix}
       & & 1 \\
       & I_{n-1} & \\
      1 & &
    \end{pmatrix}\in GL_{n+1}(\R).
  \]
  Realize left multiplication by \(X\) as a row operation and right multiplication by \(X\) a column operation, then we obtain \(XAX^{-1}\) by \begin{align*}
    A \xmapsto[R_1\leftrightarrow R_{n+1}]{X\cdot} \begin{pmatrix}
        &  & 1 \\
        & I_{n-1} & \\
      2 &  &
    \end{pmatrix} \xmapsto[C_1\leftrightarrow C_{n+1}]{\cdot X^{-1}} \begin{pmatrix}
      1 &  &   \\
        & I_{n-1} & \\
        &  & 2
    \end{pmatrix}\notin \iota(GL_n(\R)).
  \end{align*}
\end{proof}

\begin{theorem}[first isomorphism theorem]
  Let \(\f:G\to G'\) be a group homomorphism. Then the sequence \[
    \begin{tikzcd}
      G \ar[bend left=15]{rrr}{\f} \ar[->>]{r}[']{\pi} & G/\ker\f \ar{r}{\cong}[']{\wt\f} & \im\f \ar[hook]{r} & G'
    \end{tikzcd}
  \] is exact, where \(\pi\) is the canonical projection and \(\wt\f\) is the induced isomorphism given by \[
    a \xmapsto[\hspace*{2em}]{\pi} a\ker\f
    \qquad
    \wt\f(a\ker\f) = \f(a).
  \]
\end{theorem}

\begin{example}
  \(\Z/n\Z \cong \Z_n\).
\end{example}
\begin{proof}
  Consider \[
    \f:\Z\to\Z_n, \quad
    k \mapsto k\mod n.
  \]
\end{proof}

\begin{example}
  \(\R/\Z \cong U(1)\).
\end{example}
\begin{proof}
  Consider \[
    \f:\R\to U(1), \quad
    t \mapsto e^{2\pi it}.
  \]
\end{proof}

\begin{example}
  \(\R^2/\Z^2 \cong \mathbb{T}^2\), the torus.
\end{example}
\begin{proof}
  Consider \[
    \f:\R^2\to \mathbb{T}^2, \quad
    (x,y) \mapsto (e^{2\pi ix}, e^{2\pi iy}).
  \]
\end{proof}

\begin{example}
  \(\C^\times/U(1) \cong \R_{>0}\).
\end{example}
\begin{proof}
  Consider \[
    \f:\C^\times\to\R_{>0}, \quad
    z \mapsto |z|.
  \]
\end{proof}

\begin{example}
  \(GL_n(\K)/SL_n(\K) \cong \K^\times\).
\end{example}

\section{Isomorphism Theorems}

\section{Isomorphism Theorems}

\begin{definition}
  A group \(G\) is called \emph{simple} if it has no nontrivial normal subgroups.
\end{definition}

\begin{example}
  \(\Z_p\) is simple for any prime \(p\).
\end{example}

\begin{example}
  \(A_n\) is simple for \(n\geq 5\).
\end{example}

\begin{proposition}
  Any simple abelian group muust be isomorphic to \(\Z_p\) for some prime \(p\).
\end{proposition}

\begin{definition}
  A proper normal subgroup \(M\lhd G\) is called \emph{maximal} if there is no \(N\lhd G\) such that \(M\lneq N\lneq G\).
\end{definition}

\begin{proposition}
  \(M\lhd G\) is maximal if and only if \(G/M\) is simple.
\end{proposition}

\begin{definition}
  Let \(H, N\) be subgroups of \(G\). The \emph{join} \(H\vee N\) of \(H\) and \(H\) is the subgroup of \(G\) generated by \[
    HN := \{hn\mid h\in H,n\in N\}.
  \]
\end{definition}

\begin{lemma}
  If \(N\lhd G\), then \[
    H\vee N = HN = NH
  \] for any \(H < G\). If furthermore \(H\lhd G\), then \(HN\lhd G\).
\end{lemma}

\begin{theorem}
  For \(H < G\) and \(N\lhd G\), we have \[
    \frac{HN}{N} \cong \frac{H}{H\cap N}.
  \]
\end{theorem}
\begin{proof}
  Consider the map \[
    \pi: H\longrightarrow HN/N, \quad
    h \longmapsto hN.
  \]
  Clearly \(\ker(\pi) = H\cap N\) and \(\im(\pi) = HN/N\), then follow the first isomorphism theorem.
\end{proof}

\begin{example}
  Let \(m,n\in\Z\), then in \(\Z\), we have \begin{align*}
    m\Z + n\Z &= \<m,n\> = \gcd(m,n)\Z, \\
    m\Z \cap n\Z &= \<\lcm(m,n)\> = \lcm(m,n)\Z.
  \end{align*}
  The second isomorphism theorem implies that \[
    \frac{\gcd(m,n)\Z}{n\Z}
    = \frac{m\Z + n\Z}{n\Z}
    \cong \frac{m\Z}{m\Z \cap n\Z}
    = \frac{m\Z}{\lcm(m,n)\Z},
  \] which gives \[
    \gcd(m,n)\lcm(m,n) = mn.
  \]
\end{example}

\begin{theorem}
  Let \(H,K\lhd G\) and \(K < H\), then \[
    \frac{G}{H} \cong \frac{G/K}{H/K}.
  \]
\end{theorem}
\begin{proof}
  Consider the map \[
    \pi: G/K \longrightarrow \dfrac{G}{H}, \quad
    aK \longmapsto aH.
  \]
  Clearly \(\ker(\pi) = H/K\) and \(\im(\pi) = \dfrac{G}{H}\), then follow the first isomorphism theorem.
\end{proof}

\section{Free Groups and Presentations}

\begin{definition}[\emph{free group}]
  Let \(A\) be set. Consider the category \(\cal{F}^A\) with \begin{itemize}
    \item \(\Obj(\mathcal{F}^A) = \Big\{(j, G) \Bigm| G\text{ is a group and } j:A\to G\text{ is a map}\Big\}\),
    \item morphisms \((j_1,G_1)\longrightarrow (j_2,G_2)\) are commutative diagrams \[
      \begin{tikzcd}
        G_1 \ar{rr}{\f} & & G_2 \\
        &A \ar{ul}{j_1} \ar{ur}[swap]{j_2}
      \end{tikzcd}
      \text{ such that } \f \text{ is a group homomorphism.}
    \]
    A \emph{free group \(F(A)\) on \(A\)} is then the group component of an initial object in \(\cal{F}^A\). As such, it's unique up to isomorphism.
  \end{itemize}
\end{definition}

\begin{example}
  \(F(\emptyset) = \{e\}\).
\end{example}

\begin{example}
  If \(A = \{a\}\) is a singleton, then \(F(A) = \<a\> \cong \Z\).
\end{example}

If \(|A|\geq 2\), \(F(A)\) is quite complicated.

\begin{definition}
  Given a set \(A\), we define \[
    \Z^{\oplus A} = \Big\{
      f: A\to \Z : f(a)\neq 0\text{ for only finitely many } a\in A  
    \Big\}
  \]
\end{definition}

\begin{example}
  \(F(A) = \< A \mid \emptyset \>\)
\end{example}

\begin{example}
  \(\Z_n = \< 1 \mid n\>\)
\end{example}

\begin{example}
  \(D_n = \< r, s \mid r^n, s^2, (rs)^2\>\)
\end{example}

\begin{example}
  \(S_3 = \< \sgm,\tau \mid \sgm^3, \tau^2, (\sgm\tau)^2 \>\)
\end{example}